\documentclass[11pt]{article}
\usepackage[margin=1in]{geometry}
\usepackage{amsmath,amssymb}
\usepackage{booktabs}
\usepackage{natbib}
\usepackage[hidelinks]{hyperref}
% Shared notation for the differential-geometric SN Ia light-curve project.
% Input this file from every scientific document so notation stays consistent.
%
% Requires: amsmath, amssymb

% --- Number sets -------------------------------------------------------
\newcommand{\Real}{\mathbb{R}}
\newcommand{\Rn}{\mathbb{R}^{n}}
\newcommand{\Rtwo}{\mathbb{R}^{2}}
\newcommand{\Rthree}{\mathbb{R}^{3}}

% --- Differential geometry ---------------------------------------------
\newcommand{\dd}{\mathrm{d}}
\newcommand{\curve}{\gamma}                 % the light-curve curve
\newcommand{\arc}{s}                        % arclength
\newcommand{\curv}{\kappa}                  % curvature
\newcommand{\tors}{\tau}                    % torsion (n >= 3 only)
\newcommand{\frenet}[1]{\mathbf{E}_{#1}}    % Frenet frame vector
\newcommand{\tang}{\mathbf{T}}              % unit tangent
\newcommand{\norml}{\mathbf{N}}             % unit normal
\newcommand{\arclen}{L}                     % total arclength of the window

% --- Model parameters ---------------------------------------------------
\newcommand{\shape}{\theta}                 % shape latents
\newcommand{\shapeone}{\theta_{1}}
\newcommand{\shapetwo}{\theta_{2}}
\newcommand{\tscale}{\sigma}               % timing scale
\newcommand{\tmax}{t_{\mathrm{max}}}        % time origin
\newcommand{\warpone}{a_{1}}                % quadratic timing corrections
\newcommand{\warptwo}{a_{2}}
\newcommand{\dm}{\mu}                       % normalization / distance modulus
\newcommand{\colr}{c}                       % colour translation (variant only)

% --- Photometry ---------------------------------------------------------
\newcommand{\magg}{m_{g}}
\newcommand{\magr}{m_{r}}
\newcommand{\fluxg}{f_{g}}
\newcommand{\fluxr}{f_{r}}
\newcommand{\zp}{\mathrm{zp}}
\newcommand{\phase}{p}                      % rest-frame phase
\newcommand{\zred}{z}

% --- Objects and abbreviations -----------------------------------------
\newcommand{\snia}{SN\,Ia}
\newcommand{\sneia}{SNe\,Ia}
\newcommand{\ztf}{ZTF}
\newcommand{\salt}{SALT2}
\newcommand{\NN}{\mathrm{NN}}

% --- Operators ----------------------------------------------------------
\DeclareMathOperator{\SO}{SO}
\DeclareMathOperator*{\argmin}{arg\,min}
\newcommand{\deriv}[2]{\frac{\dd #1}{\dd #2}}
\newcommand{\pderiv}[2]{\frac{\partial #1}{\partial #2}}
\newcommand{\abs}[1]{\left| #1 \right|}
\newcommand{\norm}[1]{\left\| #1 \right\|}


\title{CMAGIC Distilled: \\ Empirical Evidence for Light Curves as Curves in Magnitude Space}
\date{\today}

\begin{document}
\maketitle

\begin{abstract}
A working distillation of \citet{Wang2003}, restricted to what bears on the
differential-geometric light-curve model. \salt\ supplied the parameterisation to
beat; CMAGIC supplies something different and prior --- direct empirical evidence
that the object this project treats as fundamental, the track a \snia\ traces in
magnitude space, is real, is measurable, and has exactly the invariance and
segment structure the geometric model assumes. Wang et al.\ do not fit a curve;
they fit one straight line on one post-maximum segment of the colour--magnitude
plane. But in doing so they demonstrate, on 1990s photometry, three things this
project asserts: that shape in magnitude space is insensitive to distance,
redshift, and (to first order) extinction; that \snia\ are not a one-parameter
family; and that a magnitude read off a \emph{geometric feature} of the track
standardises better than the magnitude at maximum. What is retained here is the
linear relation and its invariances, the slope and intercept as observables, the
quantified reduction in extinction sensitivity, and the honest boundary between
what a single line proves and what a full curve must still earn.
\end{abstract}

\section{The linear relation}

Wang et al.\ plot $B$ magnitude against $B-V$ colour --- the ``colour--magnitude''
or CMAG plane --- rather than either quantity against time. The central empirical
finding is that for much of the first month past $B$ maximum this track is
\emph{strikingly linear}:
\begin{equation}
  B = B_{BV} + \beta_{BV}\,(B-V),
  \label{eq:cmagic}
\end{equation}
and likewise for $B-R$ and $B-I$. The slope and intercept are constants in time for
a given supernova. Linearity holds down to the $0.01$~mag level to which the data
are sensitive (their Fig.~5): typical $\chi^{2}$ per degree of freedom is well below
one and the scatter about the line is under $0.05$~mag, so that ``higher order fits
or more complicated colour--magnitude templates are not necessary.'' The linear
regime runs roughly from day $+5$ to $+35$ after $B$ maximum for the $B$ vs $B-V$
relation (later and shorter for redder colour pairs and for fast decliners), bounded
below by a pre-maximum excursion and above by a sharp turn into the nebular phase.

The slope is defined by a ratio of rates,
\begin{equation}
  \beta_{BV} = \frac{\dot{B}}{\dot{B} - \dot{V}},
  \label{eq:slope}
\end{equation}
which is the single most important structural fact in the paper for present
purposes: \eqref{eq:slope} is the direction of the tangent to the track in the
colour--magnitude plane. CMAGIC measures a tangent direction on a segment where the
curvature is, empirically, zero.

\section{Why the shape is invariant --- the point that transfers most directly}

The morphology of the CMAG track is set by the derivatives in \eqref{eq:slope}, and
Wang et al.\ observe that this makes it insensitive to precisely the quantities this
project removes as translations and reparameterisations:

\begin{itemize}
  \item \textbf{Extinction displaces the track without deforming it.} Host and
        Galactic reddening move the entire CMAG curve along a fixed vector in the
        $(B, B-V)$ plane but leave its shape --- and hence $\beta_{BV}$ --- unchanged.
  \item \textbf{Distance and absolute luminosity} enter $B$ additively and so shift
        the track bodily, again leaving the slope untouched.
  \item \textbf{Redshift stretches the time coordinate by $1+\zred$ equally in all
        bands}, a reparameterisation that does not change the slopes in
        \eqref{eq:slope}. The shape depends on redshift and extinction only
        indirectly, through the K-correction.
\end{itemize}

This is the central design decision of the geometric model
--- that in magnitude space distance and dust act as translations, to which
curvature is blind, and time dilation and stretch act as reparameterisations, which
leave the curve pointwise unchanged --- stated and demonstrated empirically a decade
and a half earlier, in a two-band special case. Wang et al.\ reach it by direct
inspection of the data rather than by the fundamental theorem of curves, which is
exactly why it is worth citing: it is external, prior, and observational
corroboration that magnitude-space \emph{shape} is the extinction- and
distance-invariant carrier of \snia\ information.

Two honest boundaries on the corroboration. First, Wang et al.\ treat extinction as
a \emph{constant} displacement vector --- the phase-independent, zeroth-order picture.
The refinement that host reddening is a phase-dependent displacement, and therefore
deforms the curve at second order, is beyond CMAGIC and is this project's own
concern. CMAGIC supports the leading-order translation claim, not the second-order
correction. Second, their invariance to stretch is conditional: it holds only if the
$B$ and $V$ light curves stretch by the \emph{same} factor, $s_{B} \propto s_{V}$.
Where they do not, the shape changes --- which is the subject of the next section.

\section{\sneia\ are not a one-parameter family}

If \snia\ light curves were a one-parameter family --- a single stretch or $\Delta
m_{15}$ governing everything, with $s_{B} \propto s_{V}$ --- then by \eqref{eq:slope}
every CMAG track would have identical shape and $\beta_{BV}$ would be universal. Wang
et al.\ find it is not. The slope has a real dispersion beyond observational error
(K-corrected $\langle\beta_{BV}\rangle = 1.94 \pm 0.16$; without K-correction
$2.07 \pm 0.18$), it is uncorrelated with $\Delta m_{15}$, and the variations outside
the linear region --- notably a pre-maximum ``bump'' present in the brightest objects
and absent in others --- cannot be reproduced by stretching a single template
sequence. They read this as direct evidence that ``\snia\ light curves are \emph{not}
a one-parameter family,'' and note a suggestive bimodality in $\beta_{BV}$
(a branch near $1.65$ and another near $2.0$) that may separate intrinsically
distinct progenitor or burning classes.

This is the empirical warrant for carrying more than one shape degree of freedom.
The geometric model's ladder introduces $\shapeone$ and then $\shapetwo$ precisely to
test whether a second shape parameter earns its place; CMAGIC is prior evidence that
there is genuine structure for a second parameter to capture, and that at least one
axis of it --- the tangent direction on the linear segment --- is measurable and
extinction-robust.

\section{The intercept as a standardisation variable}

CMAGIC uses the \emph{intercept} of \eqref{eq:cmagic}, evaluated at a fixed reference
colour to control extrapolation error, as a distance indicator in place of the
magnitude at maximum. Wang et al.\ adopt $B_{BV0.6}$ (the fit evaluated at $B-V=0.6$,
near the sample mean colour) and a standard intercept $B_{BV} \equiv B_{BV0.6} -
1.164$~mag.

The payoff is that a magnitude defined by a geometric feature of the track
standardises \emph{better} than the peak magnitude. After a $\Delta m_{15}$
correction, the low-extinction subsample shows Hubble-residual dispersion
$\sigma \sim 0.07$~mag for the colour--magnitude intercept $M^{B}_{BV}$ against
$\sim 0.10$~mag for the maximum magnitude $M^{B}_{\max}$. On raw intercepts the
dispersion is $\approx 0.08$~mag for $(B_{\max}-V_{\max}) \le 0.05$ and $\approx
0.11$~mag for $\le 0.2$, consistent with being observational-error dominated. The
intercept is also robust for poorly observed supernovae, for which a reliable
$\Delta m_{15}$ or peak magnitude cannot be recovered.

This is the closest external precedent for the project's deliverable. The geometric
model does not standardise on peak magnitude either: its distance variable $\dm$ is
tied to the template $g$-maximum, and its Tripp-style standardisation is corrected by
latents that describe the \emph{shape of the curve}. CMAGIC establishes that reading
a distance off the geometry of the track, rather than off a single photometric peak,
is not only viable but tighter.

\section{Reduced extinction sensitivity, quantified}

Because the CMAG intercept mixes magnitude and colour, its extinction correction is
not the conventional $A_{B} = R_{B}\,E(B-V)$. Wang et al.\ show
\begin{equation}
  A_{BV} = \bigl(R_{B} - \beta_{BV}\bigr)\,E(B-V),
  \label{eq:extcorr}
\end{equation}
so with $R_{B} \approx 4.1$ and $\beta_{BV} \approx 2$ the extinction correction
applied to $B_{BV}$ is roughly \emph{half} that applied to the raw peak magnitude.
The practical consequence is that an error in $E(B-V)$ propagates into the distance
multiplied by $R_{B}-\beta_{BV} \approx 2$ rather than by $R_{B} \approx 4.1$ --- a
factor-of-two reduction in the dominant systematic for reddened objects. They further
construct an $E(B-V)$-free magnitude combination and an alternative colour-excess
estimator $\mathcal{E}(B-V) = (B_{\max}-B_{BV})/\beta_{BV}$ that tracks
$(B_{\max}-V_{\max})$ with unit slope.

This quantifies, in the two-band case, what magnitude space buys against dust: not
elimination, but a measurable reduction. It is consistent with the project's own
claim that magnitude space \emph{reduces} dust from a shape-destroying anisotropic
dilation to a small deformation about a rigid part --- reduced, not removed.

\section{What transfers and what does not}

\begin{center}
\begin{tabular}{ll}
\toprule
CMAGIC element & Status here \\
\midrule
Track lives in a colour--magnitude plane            & \textbf{shared premise} (curve in $\Rn$) \\
Linear segment with constant slope                  & \textbf{adopted} as empirical support for $\curv=0$ regions \\
Shape invariant to distance, redshift               & \textbf{adopted} --- external corroboration of the translation/reparam.\ argument \\
Shape invariant to extinction (constant vector)     & adopted at \emph{first order} only; phase-dependent dust is beyond CMAGIC \\
Not a one-parameter family                          & \textbf{adopted} --- warrant for $\shapetwo$ \\
Slope $\beta_{BV}$ as tangent direction             & \textbf{benchmark} for the model's tangent on the post-max linear phase \\
Intercept as distance variable                      & \textbf{precedent} for standardising on curve geometry, not peak \\
Reduced extinction sensitivity \eqref{eq:extcorr}   & \textbf{consistent with} the reduced-not-removed claim \\
\midrule
Fit only one line on one post-max segment           & not adopted --- the model fits the \emph{whole} curve, all epochs \\
Zero weight to data outside the linear regime       & not adopted --- curved regions carry $\curv(\arc)$, the signal \\
Linear feature is \emph{post}-maximum (day $+5$--$35$) & distinct from the model's built-in \emph{early}/pre-explosion ray \\
$B$ vs $B-V$ (needs $B$-like and $V$-like bands)     & mapped onto \ztf\ $g,r$; slope values not directly portable \\
Closed-form linear fit, transparent errors          & a caution: the neural $\curv$ is less transparent than a line \\
\bottomrule
\end{tabular}
\end{center}

Two entries deserve comment. The linear segment CMAGIC exploits is \emph{post}
maximum, whereas the straight segment built into the geometric model is the
\emph{early}, pre-explosion ray where curvature is set to zero by construction. These
are different segments of the same curve, and CMAGIC does not license the early one.
What CMAGIC licenses is the more general proposition that \snia\ tracks in magnitude
space contain extended intervals of vanishing curvature that are real and measurable
--- which makes a zero-curvature construction physically credible rather than a
convenience. It also hands the model a concrete post-maximum test: the learned
$\curv(\arc)$, mapped onto the $g,r$ analogue of the $(B,B-V)$ plane, should approach
zero over the corresponding phase interval, and the tangent direction there should
reproduce a CMAGIC-like slope. That is an independent check on the fitted geometry,
not an input to it.

The physical reading Wang et al.\ offer for the linearity --- that just after maximum
the ejecta seen at the photosphere lie in zones of more complete nuclear burning and
so are more uniform, echoed by the drop in spectropolarimetric asphericity after
maximum --- is worth keeping in view: it suggests the low-curvature regions are where
the model's shape description should be most reliable, and the high-curvature regions
near maximum and the secondary bump are where intrinsic diversity, and thus the
latents, do their work.

\section{Role of CMAGIC in this project}

CMAGIC is not a benchmark to fit against and is not used in the pipeline. Its role is
evidential. It is the prior empirical demonstration that the geometric model's
foundational choice --- represent the light curve as a curve in magnitude space and
let shape live in the geometry of that curve --- describes something real, invariant
to the right nuisances, and richer than one parameter. Where \salt\ is the model to
beat, CMAGIC is the observation that the beating is worth attempting. Its concrete
numbers ($\langle\beta_{BV}\rangle$, the intercept dispersions, the factor-of-two
extinction reduction) are targets the geometric quantities can be compared against
once fitted, on the post-maximum linear phase where the two descriptions overlap.

\bibliographystyle{plainnat}
\bibliography{refs}

\end{document}
